\documentclass[11pt]{article}
\usepackage{amsmath,amssymb,amsthm,geometry}
\geometry{margin=.9in}\setlength{\emergencystretch}{1.6em}
\newtheorem{theorem}{Theorem}\newtheorem{lemma}{Lemma}
\newcommand{\dd}{\,d}\newcommand{\E}{\mathbb E}
\title{Signed probe motion resolved against the measured weak response}
\author{AN04, E5 step 6}\date{Version 1}
\begin{document}\maketitle

\paragraph{Setting and benchmark.}
Use the actual Gaussian continuum flow with aligned balanced isotropic
critical initialization $S_0=cq^{1/\omega}$, fixed $c>0$ and
$(\rho,\zeta)$ on WIT's learned rectangle. Write $\lambda=\rho^2$,
$\omega=(1-\lambda)/2$, $L=\lambda/2$, $D(s)=e^{Ls}/(1+s)$.
Use a fixed COD witness $t_0=t_h+s_0$ chosen before $q$ so that
$e_0=e_q(t_0)<e_*=\min_J\mathcal E_{\rm coh}$ with a strict margin
for small $q$. Here $e_q=(|r_t(\xi)|^2-1)/(2q^2)$,
$r_t(z)=z-f_t(z)$, $\xi=(\sin(q\zeta),-\cos(q\zeta))$.
Set $R_0=|r_{t_0}(\xi)|$, $u=r_{t_0}(\xi)/R_0$ and
$v=(-u_2,u_1)$. Both vectors are fixed after choosing the witness,
although they depend on $q$.
For $G=f_T(\xi)-f_{t_0}(\xi)$ denote the longitudinal and transverse
changes by $B=u^TG/q^2$ and $A=v^TG/q$, respectively. The target
set of times is DRR's measured $t_{1/2}+J$, restricted to the two
response and two cluster-loss criteria; the set is empty if not
reached. Nothing below assumes a logistic weak seed or imports the
deep-family passage. Estimates stated through learning apply to any
such reached finite $T$. They do not prove that a learned time exists.

\section{Two separate signed estimates before weak learning}
\label{an04spmv1:early}
Put $H_t=f_{t,1}(\xi)/q$, $V_t=f_{t,2}(\xi)/q^2$.
For each fixed $s\ge s_0$, along any BRW extraction at $t_h+s$,
COD gives limits $H(s)=\int(\zeta-x)_+\dd\mu$ and
$V(s)=\int y(\zeta-x)_+\dd\mu$. Its exposure bound gives
$0\le H(s)\le C\sqrt{1+s}$, and
$V(s)\le C-kD(s)$ by COD's signed core bound and positive-Q1 estimate
(or by its energy bound and $H^2/2-\zeta H\ge-\zeta^2/2$).
Constants here and below may depend on the fixed family and witness.
At the witness $u_1/q\to\zeta-H(s_0)$ and $u_2\to-1$ on this
extraction. More precisely $u_1=O_{s_0}(q)$, $1+u_2=O_{s_0}(q^2)$
by the fixed-time bounds. Consequently
\begin{align*}
 A&\longrightarrow H(s)-H(s_0),\\
 B&\longrightarrow (\zeta-H(s_0))[H(s)-H(s_0)]-[V(s)-V(s_0)] .
\end{align*}
Thus the full fixed-$s$ limsup of $|A|$ is at most $C_{s_0}\sqrt{1+s}$
and the liminf of $B$ is at least $kD(s)-C_{s_0}\sqrt{1+s}$.
This proves a signed longitudinal surplus against the transverse
square in the ordered pre-weak limit. It does not evaluate the
limits at a positive weak-response hit.

\section{A regression accessible through weak learning}
Define the auxiliary weak response in the fixed probe directions by
\[
 Q=\E_{P_2}(\xi^TD)^2=\rho^2\xi_2^2+q^2,\qquad
 h(t)=\frac{\E_{P_2}[(\xi^TD)u^Tf_t(D)]}{Q}.
\]
Uniformly in time, $h(t)=\mathsf m_2(t)+O_{s_0}(q)$.
Indeed the training loss decreases along the actual gradient flow,
so $\|f_t\|_{L^2(P_2)}\le C$ by the initial loss and Gaussian moments.
Expand the two scalar products about $-e_2$, use $u_1=O_{s_0}(q)$,
$\xi_1=O(q)$, $u_2,\xi_2=-1+O_{s_0}(q^2)$ and
$Q=\rho^2+q^2+O(q^2)$. Cauchy--Schwarz gives the claim for this
order-one response comparison. In the identities below $h$ itself
is used without approximation. Since $\mathsf m_2(t_0)=o(1)$,
$\Delta h=h(T)-h(t_0)\ge1/4-o(1)>0$ at any learned time.

Use Cartesian label vectors $\Theta=(U,W)$ with
$dl=d\alpha/(2\pi)$, $|U|=|W|=\sqrt m$, $d\mu=m\,dl$.
The inner product below on paths is integrated over $dl\,dt$
from $t_0$ to $T$, with Euclidean inner product in the two vector
blocks. For a unit direction $d$, the probe gradient in this space
at each instant is
\[
 g_d=\big((d^TW)\xi\,{\bf1}_{U^T\xi>0},\
                        d(U^T\xi)_+\big).
\]
The corresponding regression gradient is
\[
 l=\left((u^TW)\frac{\E_{P_2}[(\xi^TD)D{\bf1}_{U^TD>0}]}{Q},
         \quad u\,\frac{\E_{P_2}[(\xi^TD)(U^TD)_+]}{Q}\right).
\]
These are the two physical-layer derivatives, with a.e.\ chain
rule at a probe contact. Let $F=\|l\|^2$, $\mathcal D=\|\dot\Theta\|^2$.
Then $\Delta h=\langle l,\dot\Theta\rangle$, and $\mathcal D$ is
the drop in population training loss on this interval by gradient
flow. At a learned $T$, $F>0$, $F\ge(\Delta h)^2/\mathcal D$, and
$0<\beta=\Delta h/F\le\mathcal D/\Delta h\le C$ with $C$
independent of $T,q$ for small $q$. No monotonicity of $h$ was used.

\paragraph{Gate integrals of the response component.}
Write $a=U/|U|$, $b=W/|W|$, $z=a^T\xi$, and ignore zero masses.
For $i,j\in\{u,v\}$ define the time-integrated matrix
\[
 J_{ij}=\int_{t_0}^T\int_{z>0}(i^Tb)(j^Tb)
       \frac{\E_{P_2}[(\xi^TD)^2{\bf1}_{a^TD>0}]}{Q}\,d\mu\,dt .
\]
Define the scalar direct-output contribution
\[
 \Omega=\int_{t_0}^T\int z_+\,
       \frac{\E_{P_2}[(\xi^TD)(a^TD)_+]}{Q}\,d\mu\,dt,\qquad
 I_S=\int_{t_0}^T\int d\mu\,dt.
\]
Taking the block inner products gives, exactly,
\begin{equation}\label{an04spmv1:measured}
 \langle g_u,l\rangle=J_{uu}+\Omega,\qquad
 \langle g_v,l\rangle=J_{vu},\qquad
 J\succeq0,\quad |J_{vu}|\le\sqrt{J_{uu}J_{vv}} .
\end{equation}
The last two signs are the Gram inequality. There is no sign
assignment to the mixed gate overlap $J_{vu}$.

\begin{lemma}[Both tails of the direct-output term]\label{an04spmv1:overlap}
For all the reached times in these formulas,
$-Cq^2 I_S\le\Omega\le C e^{-d_0/q^2}I_S$ for a fixed $d_0>0$.
\end{lemma}
\begin{proof}
It suffices to bound each unit donor. If $a_2\ge0$ and $z>0$,
then $a_2\le Cq$ and $z\le Cq$, since $\xi_2$ stays bounded
away from zero on its negative side. The weak-cluster activation
has $L^2$ norm at most $\rho a_2+q\le Cq$. If $a_2<0$ then
$z_+\le C(q+|a_2|)$ and its weak activation has norm
$q\| (Z-\rho |a_2|/q)_+\|_2$ for standard normal $Z$.
By Cauchy--Schwarz and the Gaussian tail, this is at most
$Cq\exp[-(\rho |a_2|/q)^2/8]$.
Multiplying by $z_+$ and $\|\xi^TD\|_2/Q$ gives $Cq^2$ in
both cases. If $z\le0$ the term vanishes. For the positive side
$\xi^TD$ has mean $\rho\xi_2\le-\rho/2$ and standard deviation
$q$, so its positive-part $L^2$ norm is exponentially small;
the other factor has bounded norm. This proves both bounds with
the actual mass and elapsed time retained.
\end{proof}

\paragraph{Signed bounds on the actual motions.}
Let $\Pi^\perp$ be the orthogonal projection off the single path
$l$ on $[t_0,T]$, set $p_B=g_u/q^2$, $p_A=g_v/q$,
$\Sigma_{ij}=\langle\Pi^\perp p_i,\Pi^\perp p_j\rangle$
for $i,j=B,A$, and
\[
 \mathcal D_\perp=\mathcal D-(\Delta h)^2/F,\qquad
 Z_i=\langle \Pi^\perp p_i,\,\dot\Theta-\beta l\rangle .
\]
All the quantities are calculated on the critical trained trajectory.
For every finite learned $T$ one has
\begin{align}
 B&=\beta(J_{uu}+\Omega)/q^2+ Z_B,&
 A&=\beta J_{vu}/q+Z_A,\label{an04spmv1:split}\\
 |d_B Z_B+d_A Z_A|&\le
 \sqrt{\mathcal D_\perp\,d^T\Sigma d}
 &&(d=(d_B,d_A)^T\in\mathbb R^2). \label{an04spmv1:ellipse}
\end{align}
Indeed $\dot\Theta-\beta l$ has squared norm $\mathcal D_\perp$;
pair it with any linear combination of the projected probe gradients.
In particular, writing $r_i=\sqrt{\mathcal D_\perp\Sigma_{ii}}$,
\begin{align*}
 \beta J_{uu}/q^2-C\beta I_S-r_B
 &\le B\le\beta(J_{uu}+C e^{-d_0/q^2}I_S)/q^2+r_B,\\
 \beta J_{vu}/q-r_A&\le A\le\beta J_{vu}/q+r_A .
\end{align*}
The response-associated input term in $B$ is favorable at a
learned endpoint. The direct-output term is paid only per unit mass
per unit time. The decomposition uses one path projection over the
whole interval, not a division by $\dot h$ or a positive-time-weight
assumption on the instantaneous weak response. The correlations $Z_i$
include all sources of motion invisible to that one integrated
regression measurement.

\section{Which part of the Gaussian training field drives the motion?}
The path projection above is a measurement identity. A separate
instantaneous decomposition identifies the source terms of the
actual gradient flow. Write $h_1(D)=D_1$, $h_2(D)=\xi^TD$,
$Q_p'=\E_{P_p} h_p^2$ and define the full-vector regressions
\[
 c_p=\E_{P_p}[h_p r_t(D)]/Q_p',\qquad
 \epsilon_p(D)=r_t(D)-c_p h_p(D).
\]
Thus $c_1=(1-\mathsf m_1,-r_{21})$ in the physical basis, with
P's convention $r_{21}=\E_{P_1}[D_1(f_2-D_2)]/(1+q^2)$.
Write $c_2=\alpha u+\delta v$. Uniformly in actual time,
\begin{equation}\label{an04spmv1:cross}
 \alpha=1-\mathsf m_2+O_{s_0}(q),\qquad
 \delta=\E_{P_2}[D_2 f_{t,1}(D)]/(\rho^2+q^2)+O_{s_0}(q).
\end{equation}
For the first formula subtract $h(t)$ from
$\E[(\xi^TD)u^TD]/Q$. For the second, $v=e_1+O(q)$ and
$\xi=-e_2+O(q)$, $\E D_1D_2=0$ on $P_2$; the uniform loss bound
controls the remaining errors. Below we use the exact regressions,
not these $O(q)$ approximations inside scaled time integrals.
In particular $\alpha\ge1/3$ from $t_0$ through the first exact
half-response hit for small $q$, without requiring monotonicity
of the trajectory. There is no such claimed lower bound after the hit.

Let $j_{ij}(t)$ and $\omega_t$ be the instantaneous integrands in
$J_{ij}$ and $\Omega$. For the strong cluster put
\begin{align*}
 d_1&=\int z_+\E_{P_1}[(a^TD)_+D_1]\,d\mu,\qquad
 M_1=\int_{z>0}bb^T\E_{P_1}[(\xi^TD)D_1{\bf1}_{a^TD>0}]\,d\mu,\\
 V_1&=(d_1I_2+M_1)c_1/2,\qquad S=\int d\mu .
\end{align*}
All lower-case kernels here are instantaneous. With
$\epsilon_p$ as defined above, direct insertion in the exact probe row
gives for $i=u,v$
\begin{equation}\label{an04spmv1:sources}
 i^T\dot f_t(\xi)=
 i^TV_1+\tfrac Q2\{\alpha j_{iu}+\delta j_{iv}
                        +\omega_t\,i^Tc_2\}+E_i,\qquad
 E_i=\tfrac12\sum_p\E_{P_p}[i^Tk_t(\xi,D)\epsilon_p(D)] .
\end{equation}
In particular the finite-noise physical weak residual, and not only
the measured path projection, gives a positive longitudinal input
contribution $Q\alpha j_{uu}/2$ before half response.
The mixed angular contribution is $\delta j_{uv}$ longitudinally
and $\alpha j_{vu}$ transversely. In the transverse equation
$\delta j_{vv}$ has the sign of the cross regression.

\paragraph{Actual bounds on the centered fields.}
Because $f_t$ is $S$-Lipschitz and $|f_t(z)|\le S|z|$,
\[
 \|\epsilon_p\|_{L^2(P_p)}\le C q(1+S).
\]
Indeed fit $h_p$ by the vector coefficient $r_t(\mu_p)/h_p(\mu_p)$
at the cluster center $\mu_1=e_1$ or $\mu_2=\rho e_2$.
The denominator is bounded away from zero, and the Lipschitz
and Gaussian deviation bounds give the displayed error for this
candidate; least squares only decreases it. On $P_1$ subtract
$(a_1)_+D_1$ from $(a^TD)_+$ in the direct part since
$\E_{P_1}D_1\epsilon_1=0$. Their $L^2$ difference is $O(q)$
by Lipschitz continuity about $e_1$. Also
$\|\xi^TD\|_{L^2(P_1)}=O(q)$. Both centered strong terms are
therefore at most $C q^2(1+S)S$ in absolute value. On $P_2$ the
direct term costs $Cq^3(1+S)S$ by the opposite-gate estimate in
Lemma~\ref{an04spmv1:overlap}. For its angular part Cauchy--Schwarz
first on data and then on labels gives
\begin{equation}\label{an04spmv1:centered}
 |E_i|\le C q^2(1+S)S+
             Cq(1+S)\sqrt{S j_{ii}},\qquad i=u,v .
\end{equation}
Here the second term retains the overlap with the actual receiver
gates; no independence of $\epsilon_2$ and those gates is used.

For example, on the portion with $\mathsf m_2\le1/2$, the pointwise
longitudinal bound simplifies to
\begin{equation}\label{an04spmv1:one}
 u^T\dot f_t(\xi)\ge
 u^TV_1+\tfrac{Q\alpha}{4}j_{uu}
             -C\delta^2 j_{vv}-C q^2(1+S)^2 S .
\end{equation}
Use $|j_{uv}|\le\sqrt{j_{uu}j_{vv}}$ and absorb both that product
times $|\delta|$ and the last term of \eqref{an04spmv1:centered}
by Young's inequality against the positive angular term. Use
$\alpha\ge1/3$, the loss bound on $c_2$, and the per-time estimate
$|\omega_t|\le C q^2 S$. The transverse component has its exact
signed center $v^TV_1+Q(\alpha j_{vu}+\delta j_{vv}+
\delta\omega_t)/2$ with the two-sided error
\eqref{an04spmv1:centered}. Integration divided by $q^2$ and $q$,
respectively, bounds the motions from the actual witness through
the hit if reached. Equations \eqref{an04spmv1:sources} and
\eqref{an04spmv1:centered} also apply after it, retaining both
signs of $\alpha$. This exhibits an unpaid penalty
$q^{-2}\int\delta^2 j_{vv}\,dt$ in the one-sided attempt, and an
unpaid signed $q^{-1}\int\delta j_{vv}\,dt$ in the transverse center.
A bound on $\mathsf m_2$ alone estimates neither; the strong
regression term and the integrated centered errors also require
budgets. The loss bound only gives $\delta=O(1)$. The next theorem tests the
smaller scale actually needed for a bounded fluctuation budget.

\section{A cross term that actually exceeds a bounded budget}
\begin{theorem}[Critical initialized flow]\label{an04spmv1:divergence}
Let $T_q$ be any sequence of reached learned times, or first weak
half-response hits. Along such a sequence with $q\downarrow0$,
\[
 \int_{t_0}^{T_q}\frac{\delta(t)_+}{q}j_{vv}(t)\,dt\to\infty,
 \qquad
 \int_{t_0}^{T_q}\frac{\delta(t)^2}{q^2}j_{vv}(t)\,dt\to\infty .
\]
In fact their lower bounds are already obtained on fixed pre-weak
intervals by taking noise small first. The first assertion concerns
the positive variation of one physical transverse source in
\eqref{an04spmv1:sources}; the second concerns the nonnegative
penalty of \eqref{an04spmv1:one}. Neither is the net probe movement.
\end{theorem}
\begin{proof}
First, at each fixed $s=t-t_h$, along BRW extraction,
\[
 \frac{\E_{P_2}[D_2 f_{t,1}(D)]}{q(\rho^2+q^2)}
       \longrightarrow \frac{K_s}{\rho},
 \qquad K_s=\int K(\rho x)\,d\mu .
\]
Indeed a bounded strong donor has $P_2$ activation divided by $q$
tending to $(Z_1+\rho x)_+$ per unit vector; its output direction
tends to $e_1$ and $D_2\to\rho$. The right endpoint strips
of initial length $\eta$ cost $O_s(\eta)$ by the physical bounds
$|U|,|W|\le C_s$, $|U_2|,|W_2|\le C_s q$ on $P_2$, and the whole
left mass is $o_s(q)$ by HBD--BRW. Compact-time convergence and
then strip exhaustion give the assertion. By COD's protected first
input moment, $K_s/\rho\ge c_3D(s)$. Since the error in
\eqref{an04spmv1:cross} is at most $C_{s_0}q$ uniformly in time,
\[
 \liminf_{q\downarrow0}\delta(t_h+s)/q\ge c_3D(s)-C_{s_0}.
\]
The error here does not depend on the later fixed $s$. In particular,
for all sufficiently large fixed $s$ this lower bound exceeds
$c_3D(s)/2$. It uses the entire true donor law.

Second, retain the polynomial factor in the core input bound.
In the limiting logistic resident law set $M=(1+e^{-s})^{-1}$,
$X_m=\int x\,d\mu$, $Y=\int y\,d\mu$. The BRW rank ceiling
integrated over $f$ gives
$(X_m)_+\le C e^{Ls}(1+s)^{-2/3}$ (use $|df|=d\pi$ on Q1,
its damped $f^{-1/3}$ integral, and the fixed-rank core). The BRW
budget bounds $Y$ above by the same order. COD's core output
equation, integrated with the logistic kernel at most one and the
nonnegative $\rho K(\rho x_a)$ source, yields
$(y_a)_-\le C e^{Ls}(1+s)^{-2/3}$. For the convolution without
decay kernel, split the integral of $e^{Lr}(1+r)^{-2/3}$ at $s/2$;
the earlier piece is exponentially smaller for large $s$, and the
later piece is bounded by $C e^{Ls}(1+s)^{-2/3}$.
At $x_a<0$ the exact bottom equation in COD therefore gives
\[
 2x_a'\ge -C e^{Ls}(1+s)^{-2/3}\Phi(\rho x_a).
\]
Choose $s_1$ so that $h(s)=s-\frac{2}{3L}\log(1+s)+d$
has $h'\ge1/2$ for $s\ge s_1$. Enlarge $d$ so $h\ge1$ there,
$x_a(s_1)>-\sqrt{h(s_1)}$ by COD's uniform lower bound, and
$C e^{-Ld}/(\rho\sqrt{2\pi})<1/2$. At the only lower face
$x_a=-\sqrt h$, the Gaussian tail bound
$\Phi(-z)\le\phi(z)/z$ makes the displayed doubled velocity
strictly greater than $-h'/\sqrt h$. Hence this lower face holds
on all finite continuations of the extracted law. Conversely
$\Phi(-z)\ge c\phi(z)/(1+z)$ for $z\ge0$ by integrating the
density from $z$ to $z+(1+z)^{-1}$. We obtain
\[
 \Phi(\rho x_a)\ge c e^{-Ls}(1+s)^{1/6}\qquad(s\ge s_1).
\]
All constants are independent of the extraction; the mass and
entrance here are precisely BRW's logistic trained limit.

At fixed $s$ the active Q4 core on a compact interior Q4 band
contributes at least $c\Phi(\rho x_a)$ to the limit inferior of
$j_{vv}(t_h+s)$. Its radial mass is bounded below independently
of $s$ by HBD clock. COD gives $x_a\le a_\rho<\zeta$; $v^Tb\to1$,
and the normalized squared weak Gaussian gate moment tends to
$\Phi(\rho x_a)$ on this band. This last limit follows at each fixed
$s$ from $D=(q Z_1,\rho+qZ_2)$ and compact tagged extraction.
The remaining contributions to $j_{vv}$ are nonnegative.

Combine the lower bounds at each sufficiently large fixed $s$.
For the two nonnegative integrands of the theorem, the respective
limit inferiors are at least $c(1+s)^{-5/6}$ and
$c e^{Ls}(1+s)^{-11/6}$. Integrate first over a fixed
$[s_2,S]$, $s_2>\max(s_0,s_1)$ sufficiently large, using Fatou's
lemma. BRW delay puts this interval before $T_q$ for sufficiently
small $q$. Finally let $S$ increase: both lower integrals diverge.
There was no angular-field approximation evaluated at $T_q$.
\end{proof}

\section{A critical-initialization obstruction to the norm estimate}
\begin{theorem}\label{an04spmv1:width}
For the fixed COD witness there is $d>0$, independent of small $q$
and of the finite learned endpoint $T$, such that
\[
 \mathcal D_\perp\ge d,\qquad
 d'{}^T\Sigma d'\ge d q^{-2}|d'|^2\quad(d'\in\mathbb R^2).
\]
Thus \eqref{an04spmv1:ellipse}, even after removing the exact
observed weak-response change and using the actual total loss drop,
has radii at least of order $1/q$ in both scaled directions.
This is a lower bound on the uncertainty of the norm estimate.
It is not a lower bound on either actual signed correlation.
\end{theorem}
\begin{proof}
Work first on the unit interval $I=[t_0,t_0+1]$.
BRW delay places this interval before every learned $T$ for small
$q$. HBD--BRW physical bounds give $S\le C$, $|U|,|W|\le C$,
$|U_2|,|W_2|\le Cq$ uniformly there. Constants in this proof may
depend on $s_0$. The instantaneous label norm of $l$ is at most
$Cq$: its input block uses $u^TW=O(q)$ and its output block
uses $\|(U^TD)_+\|_{L^2(P_2)}=O(q)$. These are all-label estimates.

The following two tagged sets are available uniformly on $I$.
A compact Q4 core band $C$ has a uniform positive lower radial
weight and an active probe margin. In the limit its common input
is $x_a\le a_\rho<\zeta$, so $(U^T\xi)_+\ge c_I q$ and
$|v^TW|\ge c_I$ for every label in $C$ at small $q$.
A second compact Q1 band $D_+$ has a uniform inactive margin
and positive lower radial weight with bounded scaled coordinates.
For this assertion one cannot reuse a fixed half-mass band at
arbitrary later times. Instead set the fixed BRW horizon to
$s_0+1$ and choose a compact Q1 band sufficiently high that
HBD outer at half mass puts both coordinates above BRW's $R_{s_0+1}$
with strict slack. Use the BRW protected barrier with entry level
$d_{\rm gate}>\sqrt2(\zeta+1)$ in its definition. The two faces
are checked by BRW on this same logistic law and finite horizon;
they give $x>\zeta+1$ throughout. Bounded charts on the band,
positive radial weights and the strict gate margins transfer to
the original $I$ by compact-time tagged extraction and the half-mass
clock. All choices precede the small-noise limit.

A definite fraction of the squared label norm of $l$ is then on
$D_+$, on which the probe gradient is zero. Indeed on that compact
band, in the notation $U=\sqrt m(\cos(qx),\sin(qx))$,
\[
 q^{-1}\E_{P_2}[(\xi^TD)(U^TD)_+]/\sqrt m
           =-\rho K(\rho x)+o(1)
\]
uniformly by the two Gaussian coordinates and Lipschitz positive
part. The output block of $l$ alone gives
$\|{\bf1}_{D_+} l(t)\|_{\rm lab}^2\ge c_I q^2$.
The total squared norm is at most $C_I q^2$ as seen above.
For any probe-gradient combination $p$ (supported on the probe-active
labels, up to the zero-gate selector), Cauchy--Schwarz on the
complement of $D_+$ therefore gives, for every coefficient $a$,
\[
 \|p(t)-a l(t)\|_{\rm lab}^2\ge c_I\|p(t)\|_{\rm lab}^2 .
\]
This is an inequality at each time, hence also holds for the
constant coefficient chosen by full-path projection.

For $p=d_B p_B+d_A p_A$, take $|d'|=1$.
On $C$, the output block has magnitude at least $c_I|d_B|/q$
by orthogonality of $u,v$. The input block of $p_A$ has magnitude
at least $c_I/q$ there and that of $p_B$ at most $C_I/q$ since
$|u^TW|\le C_Iq$. If $|d_B|$ is below a sufficiently small fixed
number, $|d_A|$ is bounded below and the input block of $p$ has
magnitude at least $c'_I/q$ by the reverse triangle inequality.
Otherwise the output block suffices. Integration over this fixed
positive-length band and $I$ yields the stated lower bound on
$\Sigma$ for every $T$.

Finally the strong response changes by a positive fixed amount on
$I$ for small $q$. Here is the required original-flow justification.
On each compact right interior the HBD Gaussian evaluations yield
$\E_{P_1}[D_1 W_1(U^TD)_+]/(1+q^2)=m+o(1)$ uniformly.
The omitted right strips have $O_I(\delta)$ cost, by the per-label
physical bound times their length; the left mass tends to zero by
HBD--BRW. The true response therefore converges at the two
endpoints to the logistic values $M(s_0)$ and $M(s_0+1)$.
Their strict difference, however small for this fixed $s_0$,
is independent of $q$. The instantaneous parameter gradient of
this response is bounded in label norm on $I$: differentiate
$\E_{P_1}[D_1 e_1^Tf_t(D)]/(1+q^2)$ and use $S\le C_I$ and
Gaussian second moments. Cauchy--Schwarz gives
$\|\dot\Theta\|_{I}\ge c_I$. Meanwhile $\|\beta l\|_I\le C_I q$,
using the upper bound on $\beta$ for every learned endpoint.
The reverse triangle inequality on $I$ bounds
$\|\dot\Theta-\beta l\|_{[t_0,T]}$ from below by $c_I/2$,
as required. Early strong training thus spends some dissipation
outside the weak-response direction even for the actual initialization.
\end{proof}

\section{Exactly which sign is still missing?}
For clarity the unmeasured terms have an explicit residual in data
coordinates. With $k_t(\xi,D)$ the full two-layer probe row of DRR
and $r_t(D)=D-f_t(D)$, form
\begin{equation}\label{an04spmv1:residual}
 W_T=\tfrac12\int_{t_0}^T\!\left\{
  \E_{P_1}[k_t r_t]+
  \E_{P_2}\!\left[k_t\left(r_t-2\beta(\xi^TD)u/Q\right)\right]
                         \right\}dt .
\end{equation}
Then $Z_B=u^TW_T/q^2$, $Z_A=v^TW_T/q$.
Indeed the actual parameter velocity is the adjoint feature gradient
applied to the full mixture residual. Subtracting $\beta l$ subtracts
the displayed weak-cluster mode with the factor two required by
the mixture. Pair against the two probe gradients. Path orthogonality
to $l$ does not imply orthogonality to a gated probe row. In
particular the $P_1$ term and the residual remaining on $P_2$ must
both be included; the latter is not assumed to vanish even when
the measured weak response is known.

Put $m_*=e_*-e_0>0$, $a_T=\beta J_{vu}/q+Z_A$, and
$D_T=1+2q^2 e_*-q^2 a_T^2$. The exact remaining inequalities
for a learned $T$ are
\begin{equation}\label{an04spmv1:missing}
 \begin{gathered}
 D_T>0,\\
 R_0-\sqrt{D_T}-\beta(J_{uu}+\Omega)
  < u^T W_T <
 R_0+\sqrt{D_T}-\beta(J_{uu}+\Omega).
 \end{gathered}
\end{equation}
Use $R_0^2=1+2q^2e_0$ and complete the square in DRR test:
$(R_0-q^2B)^2+q^2A^2<1+2q^2 e_*$. For the near side of the
permitted interval one can equivalently rationalize
$R_0-\sqrt{D_T}=q^2(a_T^2-2m_*)/(R_0+\sqrt{D_T})$.
Thus the favorable longitudinal input term has a known sign.
What remains unsigned is the longitudinal action of
\eqref{an04spmv1:residual} relative to this threshold after the
actual transverse overlap and transverse residual have been paid.
There is also no small learned-window bound here on $\Omega/q^2$
or on $J_{vv}$ integrated through the critical passage.
The physical-field shortcut that assigns a bounded total budget to
the positive transverse cross term or the squared longitudinal penalty
is false by Theorem~\ref{an04spmv1:divergence}. The orthogonal norm
bound does not deliver the needed scaled accuracy: the sufficient route requiring $r_A,r_B=O_{s_0}(1)$
fails by Theorem~\ref{an04spmv1:width}. A large favorable center
could in principle change a sign test; the width result does not
rule that out, nor does it assert that \eqref{an04spmv1:missing}
fails on the trained flow. No single bad coordinate or attribution
to a particular cohort follows from these estimates.

\paragraph{Scope.}
The actual trained family supplies the witness orientation, both
tagged bands, the strong and weak responses, and the path and
dissipation estimates. COD main and trained and BRW field, delay
and protected tail enter only on fixed resident intervals.
HBD clock, outer and physical are used on their prescribed right
ancestry or all-label physical domains. DRR projections and test
are finite-noise identities, P exactflow and learning use the
balanced Gaussian continuum with the a.e.\ probe chain rule, and
WIT supplies only the coherent benchmark. All constants may depend
on the fixed critical $c,\rho,\zeta,s_0$. No other initialization is
used to infer a sign. Existence of the critical learned window and
the strict retention inequality on it remain open.
\end{document}
